%\usepackage [ISO-8859-1]{inputenc}
\documentclass[a4paper,12pt,oneside]{amsbook}
\usepackage[T1]{fontenc}
\usepackage{listings}
\usepackage{textcomp}
\usepackage{multicol,calc}
%\usepackage{euler}
%\usepackage{garamond}
\usepackage{hyperref}

\usepackage{newcent}
%\fontfamily{fxm}\selectfont
\usepackage{fancyhdr}
\pagestyle{fancy}
% Detta paket skoter grafiken
%\usepackage[dvips]{graphicx}
\usepackage[pdftex]{graphicx}

\usepackage{amsthm}
\usepackage{amsmath}
\usepackage[swedish]{babel}
% For bortkommentering
\usepackage{verbatim}
% For att slippa indentering av Sections och SubSections
\usepackage{indentfirst}

% Har stalls paragrafindentering och avstand mellan stycken och rader mm
\setlength{\parindent}{0pt}
\setlength{\parskip}{6pt}
\setlength{\baselineskip}{12pt}
% Har definieras sidan
\setlength{\voffset}{-15mm}
\setlength{\hoffset}{-5.5mm}

\setlength{\topmargin}{0mm}
\setlength{\headheight}{10mm}
\setlength{\headsep}{8mm}
\setlength{\footskip}{16mm}

\setlength{\evensidemargin}{0mm}
\setlength{\oddsidemargin}{0mm}

\setlength{\marginparsep}{0mm}
\setlength{\marginparwidth}{0mm}

\setlength{\textwidth}{170mm}
\setlength{\textheight}{230mm}
\setlength{\paperwidth}{210mm}
\setlength{\paperheight}{297mm}
\setlength{\headwidth}{170mm}
% Huvud och Fot
%\fancypagestyle{}{
\fancyhf{}\fancyhead[c]{\small{\textit{Programmeringsolympiaden Kvalificering 2025}}}
\fancyfoot[c]{\thepage}
\renewcommand{\headrulewidth}{0.4pt}
\renewcommand{\footrulewidth}{0.4pt}
\newenvironment{lista}
{\begin{itemize}
\setlength{\parindent}{0pt}
\setlength{\itemsep}{6pt}}
{\end{itemize}}
% Uppgifter
\newcounter{probnr}
\newenvironment{tal}{%
\begin{list}
%{\textbf{\arabic{section}.\arabic{probnr}}} {\usecounter{probnr}
{\textbf{\arabic{probnr}}} {\usecounter{probnr}
\setlength{\leftmargin}{0mm}
\setlength{\rightmargin}{0mm}
\setlength{\labelwidth}{-1mm}
\setlength{\labelsep}{1mm}}
\setlength{\itemsep}{6pt}
}{\end{list}}
% Definition av avsnitt: FAKTA, EXEMPEL, PASTAENDE
\newtheorem{fakta}{Fakta}
\newtheorem{exempel}{Exempel}
\newtheorem{problem}{Problem}
%
\newtheoremstyle{test}% NAME
{20pt}      % ABOVESPACE
{10pt}      % BELOWSPACE
{\sffamily} % BODYFONT
{0pt}       % INDENT
{\scshape}  % HEADFONT
{}          % HEADPUNCT
{\newline}  % HEADSPACE
{}          % CUSTOM-HEAD-SPEC

\theoremstyle{test}
\newtheorem{program}{Program}
\newcommand{\sv}[1]{\textsc{#1}}            % Sma versaler
\newcommand{\fe}[1]{\textbf{#1}}            % FET
\newcommand{\ku}[1]{\textit{#1}}            % KURSIV
\newcommand{\cu}[1]{\texttt{#1}}            % Courier
\newcommand{\sk}[1]{\texttt{#1}}            % Courier
\newcommand{\rubrik}[1]{\begin{center}\sf\huge{#1}\normalsize\rm\end{center}}
\begin{document}
%\DeclareGraphicsExtensions{.jpg,.pdf,.mps,.png,.eps}



\thispagestyle{fancy}
\lstset{basicstyle=\ttfamily,
  breaklines=true}
% ----------------------------------------------------

\begin{center}
\Huge{Programmeringsolympiaden 2025}
\end{center}
\vspace{-1.2cm} 
\specialsection*{Competition rules for skolkvalet}
\begin{lista}
\item The competition takes place on a date determined by the school under {\bf four hours.
    No extra time is given for lunch or breaks.} Before the competition starts, the student
    shall decide together with the teacher whether they will use their own or a school computer.
    In both of these cases, the student must write skolkvalet in an agreed-upon location
    in the school.
    
\item The competition consists of seven tasks that shall each be solved by using a computer
  program in your programming language of choice.
\item {\bf Input can be read into the program in any way}, e.g., through the program interacting
with the user (as in the sample runs in the tasks), being entered into a graphical interface,
or data files being sent to {\em standard input}. Just agree with your teacher on how the program will be tested.
\item Your solutions will be tested with prepared input data. Each task is typically tested
with 5 test cases, each awarding 1 point if your program outputs the correct answer within
an execution time limit of {\bf 3 seconds}. No input validation is required, as it adheres
to the specifications in the task.
\item There are often different constraints for the various test cases, such as the size of
the input or other limitations. These are specified in the task. {\bf Note that the difficulty
level of a task can vary significantly depending on these differences. Therefore, it might
be easier to score partial points on a seemingly difficult task than to score full points
on a seemingly easier task.} The information about partial scoring is thus extremely
important for planning your competition.
\item The grading is performed on the same or an equivalent computer as where the code was written.
Changes to the source code are not allowed after the competition. If the program cannot be
compiled, the task will be awarded 0 points.

\item If any of the following happens, 0 points are awarded on the particular {\em testcase} being tested,
  but the program shall still be tested on the rest of the testcases.
\begin{lista}
\item The execution time of the program exceeds 3 seconds \vspace{-0.2cm}
\item Execution error (run time error) \vspace{-0.2cm}
\item Wrong answer \vspace{-0.2cm}
\end{lista}

\item Participation is individual, which means, among other things, that no exchange
of ideas or files is allowed during the competition.  
\item Permitted resources: Any written material, material installed on the computer,
and material available on the internet. 
It is \ku{not} allowed to actively communicate on the internet (e.g., chatting or asking questions in a forum);
only searching for information is permitted.  
{\bf It is also not allowed to use generative AI such as ChatGPT or GitHub Copilot.}  
Calculators are allowed.  
\item Submissions must be made in the form of source code files (e.g., uppg1...uppg7 with appropriate
file extensions) uploaded to the agreed location. Other files will not be considered. Make sure to
submit the correct version of your program.  

\end{lista}

%De högst placerade i kvalet går vidare till finalen i Göteborg där landslagsplatser till BOI i Polen och IOI i Bolivia står på spel.
\begin{center}
\Large Good luck!
\end{center}
%\end{document}

\newpage
\specialsection*{Task 1 -- The Triangle Factory}

Tristian works at a triangle factory.  
His job is to classify different types of triangles produced in the factory.  
Tristian now asks you to write a program that can replace him.  

Tristian will give you three positive integers, $a$, $b$, and $c$.  
These are the three angles of a triangle, given in degrees.  

Your program should do the following:  
\begin{itemize}
  \item If the triangle is obtuse, the program should print ``Trubbig Triangel''.
  \item If the triangle is acute, the program should print ``Spetsig Triangel''.
  \item If the triangle is right-angled, the program should print ``Ratvinklig Triangel''.  
  (Since it can be tricky to write the letter "ä," both "Right Triangle" and "Ratvinklig Triangel" are acceptable answers.)  
\end{itemize}  

To remind you of the definitions of obtuse, acute, and right-angled triangles:  
\begin{itemize}
  \item A triangle is obtuse if any of its angles are greater than $90$ degrees.  
  \item A triangle is acute if all its angles are less than $90$ degrees.  
  \item A triangle is right-angled if one of its angles is exactly $90$ degrees.  
\end{itemize}  

Your program should read the three angles $a$, $b$, $c$.  

It is guaranteed that $(1 \leq a, b, c < 180)$ holds. It is also guaranteed that the input forms a valid triangle, meaning $a + b + c = 180$.  

The program should then print whether the triangle formed by the angles $a$, $b$, and $c$ is acute, obtuse, or right-angled.  

\textbf{Scoring:} The program will be run on 5 test cases, each worth $1$ point.

\vspace{1cm}
\fe{Körningsexempel 1}
\begin{verbatim}
Vinkel 1 ? 100
Vinkel 2 ? 40
Vinkel 3 ? 40

Svar: Trubbig Triangel
\end{verbatim}
Se nästa sida för fler körningsexempel.

\newpage

\fe{Körningsexempel 2}
\begin{verbatim}
Vinkel 1 ? 60
Vinkel 2 ? 40
Vinkel 3 ? 80

Svar: Spetsig Triangel
\end{verbatim}

\fe{Körningsexempel 3}
\begin{verbatim}
Vinkel 1 ? 30
Vinkel 2 ? 60
Vinkel 3 ? 90

Svar: Ratvinklig Triangel
\end{verbatim}


\newpage
\specialsection*{Uppgift 2 -- Truls trubbel}


%\begin{center}
%  \includegraphics[width=12cm]{sifferkrypto.pdf} \\ ~\\
%  \ku{Exempel på missuppfattning som kan ske med Alice nya chiffer.}
%\end{center}

Truls och Henry spelar pingis mot varandra.

Både Truls och Henry ogillar starkt matematik.
Mer än allting annat ogillar de att räkna på antalet poäng som vardera person har under tävlingen.
För att förenkla poängräkningen har de kommit på en egen lösning:

Istället för att ständigt hålla koll på hur många poäng de har, skriver de 
ett ”\texttt{T}” om Truls får poäng och ett ”\texttt{H}” om Henry
får poäng, direkt efter varje vunnen boll.

När Truls och Henry spelar pingis fungerar poängsystemet enligt dessa tre regler:
\begin{enumerate}
  \item Vid varje boll delas ett poäng ut.
  \item En spelare vinner en match om båda av de två följande villkoren är uppfyllda:
  \begin{itemize}
    \item Spelaren leder med åtminstone 2 poäng över sin motståndare.
    \item Spelaren har minst 11 poäng.
  \end{itemize}
  Till exempel, om Truls har 11 poäng och Henry har 9 poäng, så utses Truls till vinnare. Men om Truls har 11 poäng och Henry har 10 poäng finns det ännu ingen vinnare.

  \item När en spelare har vunnit enligt ovanstående kriterier, så nollställs poängen för båda spelarna och en ny match startar.
\end{enumerate}

Nu har Truls och Henry spelat väldigt länge, och båda vill veta hur poängen ser ut.
De ber de dig att skriva ett program som
räknar ut den aktuella poängställningen åt dem, baserat på deras anteckningar.

Ditt program ska först läsa in poängsträngen, som består av flera bokstäver på en rad. Antalet bokstäver är mellan 1 och 30. 
Varje bokstav kommer antingen vara ”\texttt{T}” eller ”\texttt{H}”.

Programmet ska sedan skriva ut poängen i den nuvarande matchen som spelas, ”$T$-$H$”, där $T$ är antalet poäng som Truls har, och $H$ är antalet poäng som Henry har.

\fe{Poängsättning:}\\
De fem testfallen som vardera ger $1$ poäng har följande begränsningar:
\begin{itemize}
  \item[(1-2)] $\;\;\;\;$En match avgörs alltid så fort någon får 11 poäng.
  \item[(3-5)] $\;\;\;\; $Inga särskilda begränsningar.
\end{itemize}

\vspace{1cm}

\fe{Körningsexempel 1}
\begin{verbatim}
Anteckningarna ? TTTHHHTHH

Svar: 4-5
\end{verbatim}

Se nästa sida för fler körningsexempel.


\newpage

\fe{Körningsexempel 2}
\begin{verbatim}
Anteckningarna ? TTTTTTTTTTT

Svar: 0-0
\end{verbatim}

\fe{Förklaring:} Eftersom Truls vann matchen när poängställningen blev \texttt{11-0}, så nollställdes poängen till \texttt{0-0}.

\vspace{0.5cm}

\fe{Körningsexempel 3}
\begin{verbatim}
Anteckningarna ? THTHTHTHTHTHTHTHTHTHTHTH

Svar: 12-12
\end{verbatim}
\fe{Förklaring:} Även om båda har fler än 11 poäng leder ingen med 2 poäng. Därför har ingen vinnare utsetts ännu.


\fe{Körningsexempel 4}
\begin{verbatim}
Anteckningarna ? THTHTHTHTHTHTHTHTHTHTHTHHHH

Svar: 0-1
\end{verbatim}
\fe{Förklaring:} Detta spel är en fortsättning på körningsexempel 3 där Henry vinner matchen efter ytterligare två vunna bollar. 
Poängen nollställs, och Henry vinner därefter en till boll. Därför står det \texttt{0-1}.


\newpage
\specialsection*{Uppgift 3 -- Virus}

Boschua är inte så bra på cybersäkerhet. Förra månaden stängde han av brandväggen för att den
inte lät honom ladda ner kattbilder. Tyvärr har detta lett till att han fått ett virus. 
Viruset gjorde något väldigt konstigt: det lade till en massa bokstäver i alla filnamn.

Exempelvis kan det ha ändrat filnamnet \texttt{katt.png} till \texttt{k\textbf{atteg}att\textbf{o}.p\textbf{i}ng}.
Viruset kan alltså lägga till bokstäver, men aldrig ta bort bokstäver.

Han letar nu  efter filen som heter $F$, men han har jättemånga filer. Kan du skriva ett program som läser in filnamnet $F$ och namnet på en fil $H$
på hans dator, och avgör om $F$ kan ha gjorts om till $H$ av viruset? Det vill säga, är det möjligt att
viruset tog filnamnet $F$ och lade till bokstäver så att filen nu heter $H$?

Skriv ett program som läser in $F$ och $H$. Låt $|F|$ vara antalet bokstäver i $F$, och detsamma för $|H|$.
Det gäller att $1 \le |F| \le |H| \le 32$. $F$ och $H$ innehåller endast bokstäver \texttt{a}-\texttt{z} och punkt.
Det är garanterat att $F$ innehåller exakt en punkt. Viruset kan lägga till punkter och vilken bokstav från \texttt{a}-\texttt{z} som helst.

Ditt program ska sedan skriva ut ”Ja” om $H$ kan vara filen Boschua söker, annars ”Nej”.


\vspace{0.5cm}
\fe{Poängsättning:}\\
I detta problem finns det fem \textit{par} av testfall. I varje par är svaret på det ena testfallet Ja och det andra Nej.
Detta för att förhindra att du får poäng om ditt program alltid skriver ut ”Ja” eller alltid ”Nej”.

De fem paren av testfall som vardera ger $1$ poäng har följande begränsningar utöver de som redan nämnts:
\begin{itemize}
    \item[(Par 1)] $\;\;\;\;$Varje bokstav (och punkt) dyker upp som mest en gång i vardera sträng.    
    \item[(Par 2)] $\;\;\;\;|H| = |F|+1$ 
    \item[(Par 3)] $\;\;\;\;|H| = |F|+2$ 
    \item[(Par 4-5)] $\;\;\;\; $Inga särskilda begränsningar.
\end{itemize}

\vspace{0.5cm}
\fe{Körningsexempel 1}
\begin{verbatim}
F ? katt.png
H ? kattegatto.ping

Svar: Ja
\end{verbatim}
\fe{Förklaring:} Viruset kan ha gjort följande förändring för göra om $F$ till $H$: \texttt{katt.png} $\Rightarrow$ \texttt{ka\textbf{ttega}tt\textbf{o}.p\textbf{i}ng}.

Se nästa sida för fler körningsexempel.
\newpage

\fe{Körningsexempel 2}
\begin{verbatim}
F ? antivirus.exe
H ? aaantivisiiruis.egxe

Svar: Ja
\end{verbatim}
\fe{Förklaring:} Viruset kan ha gjort följande förändring för att göra om $F$ till $H$: \\
\texttt{antivirus.exe} $\Rightarrow$ \texttt{a\textbf{aa}ntiv\textbf{is}i\textbf{i}ru\textbf{i}s.e\textbf{g}xe}.

\fe{Körningsexempel 3}
\begin{verbatim}
F ? kattenmilly.png
H ? kaaatttenmilly.pgn

Svar: Nej
\end{verbatim}
\fe{Förklaring:} $F$ och $H$ skiljer sig vid \texttt{png} och \texttt{pgn}. Eftersom viruset inte kan ändra ordningen på bokstäver
kan det inte ha ändrat $F$ till $H$. 



\newpage
\specialsection*{Uppgift 4 -- Siffrids siffersumma}

Siffrid älskar att leka med tal! Just nu sitter hon och leker med heltalet $N$.

Hon undrar hur man kan skapa det minsta och största talet som har samma siffersumma och antal siffror som $N$.
Kan du hjälpa henne?
Siffersumman av ett tal definieras som summan av alla siffror i talet. 

\noindent
Till exempel är siffersumman av talet $1234$: 
$$
1 + 2 + 3 + 4 = 10,
$$

och siffersumman av talet $220$:
$$
2 + 2 + 0 = 4.
$$

Skriv ett program som först läser in heltalet $N$ ($1 \leq N \leq 10^9$), talet som Siffrid leker runt med.

Programmet ska sedan skriva ut det minsta och det största talet som har samma siffersumma och antal siffror som talet $N$.

\vspace{0.5cm}
\fe{Poängsättning:}\\
De fem testfallen som vardera ger $1$ poäng har följande begränsningar utöver de som redan nämnts:
\begin{itemize}
    \item[(1-2)] $\;\;\;\;$ $N \leq 10^5$    
    \item[(3-5)] $\;\;\;\;$ Inga särskilda begränsningar.
\end{itemize}

\vspace{0.5cm}
\fe{Körningsexempel 1}\\
\vspace{-0.3cm}
\begin{verbatim}
N ? 101

Svar: 101 200
\end{verbatim}
\fe{Förklaring:} Vi kan inte få ett tal mindre än 101. Tal som 11 och 2 har samma siffersumma som $N$ och är mindre än 101, men de har inte
lika många siffror som $N$.

\vspace{0.3cm}
\fe{Körningsexempel 2}\\
\vspace{-0.3cm}
\begin{verbatim}
N ? 878

Svar: 599 995
\end{verbatim}
\vspace{0.3cm}

\fe{Körningsexempel 3}\\
\vspace{-0.3cm}
\begin{verbatim}
N ? 1337

Svar: 1049 9500
\end{verbatim}


\newpage
\specialsection*{Uppgift 5 -- Miniräknaren}

Det går rykten om att skolans miniräknare kan trimmas med en hemlig kod.
Du har hört att om du kan få miniräknaren att visa den hemliga kålnami-koden på displayen och sedan
trycker på likhetstecknet, så kommer miniräknaren att låsa upp flera
stycken hemliga matematikfunktioner.

Din kompis har lyckats lista ut vad kålnami-koden är, men hen berättar för dig
att du inte får gå tillväga hur som helst för att skriva in koden i miniräknaren.
Från början visar miniräknaren $0$. Sedan får du utföra följande operationer:
\begin{itemize}
  \item Multiplicera talet som visas med talet $M$.
  \item Addera ett valfritt tal $k$ till talet som visas, där $k$ uppfyller $0 \leq k \leq M-1$.
\end{itemize}

Som tur är har din kompis även lyckats lista ut vad talet $M$ är. 
Denna komplicerade process kommer endast fungera om du gör så få operationer
som möjligt. Därför vill du nu ta reda på hur många operationer som krävs,
givet kålnami-koden och talet $M$.

Skriv ett program som först läser in heltalet $N$ ($1 \le N \le 10^9$), som utgör kålnami-koden.
Det ska sedan läsa in siffran $M$ ($2 \leq M \leq 9$), som beskrivits ovan.

Programmet ska sedan skriva ut ett heltal: det minsta antalet operationer som krävs för att få miniräknaren att visa
kålnami-koden, $N$.

\fe{Poängsättning:}\\
De fem testfallen som vardera ger $1$ poäng har följande begränsningar utöver de som redan nämnts:
\begin{itemize}
    \item[(1)] $\;\;\;\;M=3, N \leq 10$ 
    \item[(2)] $\;\;\;\;M=2$ 
    \item[(3)] $\;\;\;\;N \leq 10^5$
    \item[(4-5)] $\;\;\;\; $Inga särskilda begränsningar.
\end{itemize}


\fe{Körningsexempel 1}
\begin{verbatim}
N ? 4
M ? 2

Svar: 3
\end{verbatim}

\fe{Förklaring:}
Ett sätt att skriva in koden med minsta antalet operationer är att göra som följande:
\begin{enumerate}
  \item $+1$
  \item $\times 2$
  \item $\times 2$
\end{enumerate}

Se nästa sida för ett till körningsexempel.

\newpage

\fe{Körningsexempel 2}
\begin{verbatim}
N ? 33
M ? 3

Svar: 5
\end{verbatim}
\fe{Förklaring:}
En optimal lösning är att utföra följande operationer:
\begin{enumerate}
  \item $+1$
  \item $\times 3$
  \item $\times 3$
  \item $+2$
  \item $\times 3$
\end{enumerate}


\newpage
\specialsection*{Uppgift 6 -- Trollkarlen Theodor}

En mindre armé med $N$ monster befinner sig precis utanför Tästerås! Nu är det upp
till Trollkarlen Theodor att besegra den och rädda staden. Varje monster har ett visst
antal liv. För att skada monster kan
Theodor avfyra magiska explosioner. Varje gång han avfyrar en explosion siktar han på ett
särskilt monster. Sedan händer följande:
\begin{itemize}
  \item Alla monster förlorar $A$ liv av explosionen.
  \item Monstret som Theodor siktade på förlorar ytterligare $S$ liv, 
  eftersom det befinner sig i mitten av explosionen. Alltså förlorar monstret han siktar
  på $S+A$ liv sammanlagt av just den explosionen.
\end{itemize}
Han kan fritt välja vilket monster han ska sikta på varje gång han avfyrar en explosion.
Eftersom det är jobbigt att avfyra explosioner vill han veta det minsta antalet
explosioner han behöver skjuta iväg för att besegra alla monster om han siktar optimalt.

Skriv ett program som först läser in heltalen $N$ ($1 \leq N \leq 10$), $S$ ($1 \leq S \leq 10^9$) och $A$ ($0 \leq A \leq 10^9$), 
vars betydelse beskrivs ovan. Därefter ska programmet läsa in $N$ heltal $h_1, h_2, \ldots, h_N$ ($1 \leq h_i \leq 10^9$), som
beskriver antalet liv som vardera monster har.

Programmet ska skriva ut det minsta antalet explosioner som Theodor behöver avfyra för att besegra
alla monster. Ett monster är besegrat om det har 0 eller färre liv.

\fe{Poängsättning:}\\
De fem testfallen som vardera ger $1$ poäng har följande begränsningar utöver de som redan nämnts:
\begin{itemize}
    \item[(1)] $\;\;\;\; A=0$
    \item[(2-3)] $\;\;\;\; h_i \leq 100$
    \item[(4-5)] $\;\;\;\; $Inga särskilda begränsningar.
\end{itemize}


\fe{Körningsexempel 1}
\begin{verbatim}
N, S, A ? 3 2 1
Monstrens liv ? 7 2 3

Svar: 3
\end{verbatim}
\fe{Förklaring:} Ett optimalt sätt att avfyra explosioner är följande:
\begin{enumerate}
  \item Avfyra en explosion på monstret med 7 liv. Monstren har nu \verb+4,1,2+ liv kvar.
  \item Avfyra en explosion på monstret med 4 liv kvar. Monstret med 1 liv besegras. De resterande monstren har \verb+1,1+ liv kvar. 
  \item Avfyra en explosion på vilket som helst av monstren. Nu har alla monster blivit besegrade och Theodor är klar.
\end{enumerate}


\fe{Körningsexempel 2}
\begin{verbatim}
N, S, A ? 4 3 2
Monstrens liv ? 7 8 9 10

Svar: 4  
\end{verbatim}


\specialsection*{Uppgift 7 -- Theodor, Trollkarlen}


Detta problem bygger på problemet \textit{Trollkarlen Theodor}.

Ondskefulla Olle, som är ytterst ansvarig för all ondska, har tröttnat på att Theodor alltid förstör allting.
Nu har Olle kommit på en ny ondskefull plan, och för att sätta den i verket måste han uppehålla Theodor.
Till sitt förfogande har han $N$ monster, och eftersom han är så enormt mäktig kan han ge dessa monster olika antal liv.
Olle vill nu räkna antalet sätt han kan ge olika mängder liv till sina olika monster, 
så att Theodor behöver avfyra exakt $K$ magiska explosioner för att besegra alla monster. Enda kravet på antalet liv
för ett monster är att det ska vara större än $0$.

Precis som i problemet \textit{Trollkarlen Theodor} skadar Theodor monster genom att sikta på ett särskilt monster och avfyra
en explosion. Monstret han siktar på förlorar då $S$ liv, och sedan förlorar alla monster (även det han siktade på) $A$ liv var. 
Ett monster är besegrat om det har 0 eller färre liv.
Theodor väljer hur han siktar på monster så att han minimerar antalet explosioner han behöver avfyra.

Skriv ett program som läser in heltalen $N, K, S, A$ ($1 \leq N,K,S \leq 2\cdot10^5$, $0 \leq A \leq 2\cdot10^5$), som är de fyra talen som beskrivits ovan.

Programmet ska sedan skriva ut antalet sätt Olle kan ge liv till sina monster så att Theodor behöver avfyra exakt $K$ explosioner för att besegra alla.
Eftersom detta tal kan bli väldigt stort ska det skrivas ut \textit{modulo} $10^9+7$. Om du bara vill lösa
första testfallet behöver du inte tänka på modulo eftersom svaret är litet.


\fe{Poängsättning:}\\
De fem testfallen som vardera ger $1$ poäng har följande begränsningar utöver de som redan nämnts:
\begin{itemize}
    \item[(1)] $\;\;\;\; A = 0$, $N, K, S \leq 5$
    \item[(2)] $\;\;\;\; A = 0$, $S \leq 10$, $N, K \leq 70$
    \item[(3)] $\;\;\;\; N, K, S, A \leq 150$
    \item[(4)] $\;\;\;\; N, K, S, A \leq 1000$
    \item[(5)] $\;\;\;\;$Inga särskilda begränsningar.
\end{itemize}


\fe{Körningsexempel 1}
\begin{verbatim}
N, K, S, A ? 1 1 5 0

Svar: 5
\end{verbatim}
\fe{Förklaring:} Det finns ett monster, och de möjliga liv det kan ha så att Theodor besegrar det med exakt en explosion är \verb+1, 2, 3, 4, 5+.

Se nästa sida för fler körningsexempel.

\newpage

\fe{Körningsexempel 2}
\begin{verbatim}
N, K, S, A ? 2 2 1 1

Svar: 10
\end{verbatim}
\fe{Förklaring:} De 10 möjliga fördelningarna av monsterliv så att Theodor besegrar båda med exakt två explosioner är:
\begin{verbatim}
  1 3
  1 4
  2 2
  2 3
  2 4
  3 1
  3 2
  3 3
  4 1
  4 2
\end{verbatim}

\vspace{1cm}
\fe{Körningsexempel 3}
\begin{verbatim}
N, K, S, A ? 10 20 30 40

Svar: 435462573
\end{verbatim}

\end{document}
